\chapter{Wann gelingt ein Vorhaben?}

Helden wollen ständig etwas tun: über einen Bach springen, ein Schloss öffnen, einen Streit schlichten, eine Spur finden oder sich an einem schlafenden Troll vorbeischleichen.

Manchmal ist sofort klar, was passiert. Dann erzählt ihr einfach weiter. Manchmal ist aber unklar, ob ein Vorhaben gelingt. Dann wird gewürfelt. Das nennt man eine Probe.

Eine Probe beantwortet die Frage: Gelingt dem Helden, was er versucht?

\section{Erst erzählen, dann würfeln}

Bevor gewürfelt wird, beschreibt der Spieler, was sein Held tun möchte.

Zum Beispiel:
\begin{itemize}
  \item „Ich springe über den Bach.“
  \item „Ich schleiche an der Wache vorbei.“
  \item „Ich versuche, den verletzten Wolf zu beruhigen.“
  \item „Ich stemme die alte Tür auf.“
  \item „Ich suche nach Spuren im Matsch.“
\end{itemize}

Die Spielleitung hört zu und entscheidet dann, ob überhaupt gewürfelt werden muss.

Nicht jede Handlung braucht eine Probe. Eine Probe wird nur dann gebraucht, wenn der Ausgang wirklich unklar ist und das Ergebnis die Geschichte verändert.

\subsection{Wann wird gewürfelt?}

Gewürfelt wird, wenn drei Dinge gleichzeitig zutreffen:
\begin{itemize}
  \item Der Held versucht etwas, das grundsätzlich möglich ist.
  \item Es ist unklar, ob es gelingt.
  \item Erfolg oder Fehlschlag machen einen Unterschied.
\end{itemize}

Wenn ein Held etwas Schwieriges, Gefährliches oder Spannendes versucht, ist eine Probe oft sinnvoll.

Beispiele:
\begin{itemize}
  \item Ein Held will über eine tiefe Schlucht springen.
  \item Ein Held will sich unbemerkt an Räubern vorbeischleichen.
  \item Ein Held will ein altes Rätsel verstehen.
  \item Ein Held will einen wütenden Dorfbewohner beruhigen.
  \item Ein Held will in der Dunkelheit eine versteckte Spur finden.
\end{itemize}

In solchen Situationen entscheidet die Probe, wie es weitergeht.

\subsection{Wann wird nicht gewürfelt?}

Nicht gewürfelt wird, wenn das Ergebnis bereits klar ist.

Ein Vorhaben kann einfach gelingen, wenn es leicht ist und nichts dagegen spricht.

\textbf{Beispiel:} Ein Held möchte eine normale Tür öffnen. Die Tür ist nicht abgeschlossen, nicht verklemmt und niemand hält sie zu. Dann gelingt das einfach.

Ein Vorhaben kann auch unmöglich sein. Dann wird ebenfalls nicht gewürfelt.

\textbf{Beispiel:} Ein Held möchte ohne Magie über einen ganzen Berg springen. Das ist nicht nur schwierig, sondern unmöglich. Dafür wird keine Probe gemacht.

Auch wenn ein Fehlschlag nichts verändern würde, braucht ihr keine Probe.

\textbf{Beispiel:} Ein Held sucht in einem leeren Raum nach einer geheimen Tür. Wenn dort wirklich keine geheime Tür ist, muss nicht gewürfelt werden. Die Spielleitung kann einfach sagen, dass der Held nichts findet.

Würfelt nur, wenn das Ergebnis spannend ist.

\section{Was ist eine Probe?}

Eine Probe ist ein Würfelwurf, mit dem geprüft wird, ob ein Vorhaben gelingt.

Für eine Probe nimmt der Spieler mehrere Würfel. Danach werden alle Würfel geworfen. Auf manchen Würfelseiten sind Sterne zu sehen.

Nach dem Wurf werden alle Sterne zusammengezählt. Wenn der Held mindestens so viele Sterne würfelt, wie für das Vorhaben nötig sind, gelingt die Probe. Wenn er zu wenige Sterne würfelt, gelingt die Probe nicht oder nur teilweise.

\subsection{Welche Würfel werden verwendet?}

Für eine Probe werden bis zu drei Arten von Würfeln verwendet:
\begin{itemize}
  \item der Heldenwürfel,
  \item grüne Attributswürfel,
  \item blaue Fertigkeitswürfel.
\end{itemize}

Jede Probe enthält den Heldenwürfel.

Dazu kommen grüne Attributswürfel und blaue Fertigkeitswürfel, wenn sie zur Handlung passen.

\subsection{Welche Fertigkeit passt?}

Wenn ein Held etwas versucht, entscheidet die Spielleitung gemeinsam mit dem Spieler, welche Fertigkeit am besten passt. Dabei hilft die Beschreibung der Handlung.

Beispiele:
\begin{itemize}
  \item Ein Held greift mit einem Schwert an: Schlagen
  \item Ein Held schießt mit einem Bogen: Schießen
  \item Ein Held klettert eine Mauer hoch: Turnen
  \item Ein Held schleicht an einer Wache vorbei: Schleichen
  \item Ein Held sucht eine versteckte Spur: Suchen
  \item Ein Held löst ein Rätsel: Nachdenken
  \item Ein Held versorgt eine Wunde: Heilen
  \item Ein Held will jemanden überzeugen: Überzeugen
\end{itemize}

Jede Fertigkeit gehört zu zwei Attributen. Diese beiden Attribute zeigen, welche natürlichen Eigenschaften für die Fertigkeit wichtig sind.

Für eine Probe nimmt der Spieler:
\begin{itemize}
  \item den Heldenwürfel,
  \item die grünen Würfel der beiden passenden Attribute,
  \item die blauen Würfel der passenden Fertigkeit.
\end{itemize}

\subsection{Wenn keine Fertigkeit passt}

Manchmal versucht ein Held etwas, das zu keiner Fertigkeit richtig gut passt. In dem Fall kann die Spielleitung eine Attributsprobe verlangen.

Bei einer Attributsprobe verwendet die Spielleitung nach eigenem Ermessen entweder zwei passende Attribute oder ein einzelnes besonders wichtiges Attribut.

Wenn zwei Attribute kombiniert werden, nimmt der Spieler:
\begin{itemize}
  \item den Heldenwürfel,
  \item die grünen Würfel des ersten passenden Attributs,
  \item die grünen Würfel des zweiten passenden Attributs.
\end{itemize}

Wenn nur ein einzelnes Attribut verwendet wird, nimmt der Spieler:
\begin{itemize}
  \item den Heldenwürfel,
  \item die grünen Würfel des passenden Attributs doppelt.
\end{itemize}

Für jeden grünen Würfel dieses Attributs nimmt der Spieler in diesem Fall also zwei grüne Würfel.

Beispiele:
\begin{itemize}
  \item Ein Held will einen schweren Stein anheben.
  \item Ein Held will eisige Kälte aushalten.
  \item Ein Held will trotz großer Angst weitergehen.
  \item Ein Held will sich an einem Seil über eine Schlucht hangeln.
\end{itemize}

Attributsproben sollten die Ausnahme bleiben. Meistens ist es besser, eine passende Fertigkeit zu finden.

\subsection{Der Ablauf einer Probe}

Eine Probe läuft immer in derselben Reihenfolge ab.
\begin{enumerate}
  \item Der Spieler beschreibt, was sein Held tun möchte.
  \item Die Spielleitung entscheidet, ob eine Probe nötig ist.
  \item Die Spielleitung legt fest, welche Fertigkeit oder welches Attribut passt.
  \item Die Spielleitung nennt die Schwierigkeit.
  \item Der Spieler nimmt die passenden Würfel.
  \item Der Spieler würfelt.
  \item Alle Sterne werden gezählt.
  \item Die Spielleitung beschreibt, was passiert.
\end{enumerate}

Wichtig ist: Erst wird beschrieben, dann wird gewürfelt.

Die Würfel entscheiden nicht, was ein Held versucht. Das entscheidet der Spieler. Die Würfel entscheiden nur, ob es gelingt.

\subsection{Wie schwierig ist das Vorhaben?}

Nicht jedes Vorhaben ist gleich schwer.

Die Spielleitung legt fest, wie viele Sterne nötig sind. Je schwieriger ein Vorhaben ist, desto mehr Sterne braucht der Held.

\begin{center}
\begin{tabular}{ll}
\textbf{Schwierigkeit} & \textbf{Benötigte Sterne} \\
Leicht & 1 \\
Normal & 2 \\
Schwierig & 3 \\
Sehr schwierig & 4 \\
Heldentat & 5 \\
\end{tabular}
\end{center}

Eine leichte Probe sollte meistens gelingen. Eine normale Probe ist eine echte Herausforderung. Eine schwierige Probe braucht gute Werte oder etwas Glück. Eine sehr schwierige Probe ist nur für Helden geeignet, die in diesem Bereich wirklich stark sind. Eine Heldentat ist etwas Besonderes.

Die Spielleitung sollte die Schwierigkeit so wählen, dass sie zur Situation passt.

\subsection{Ergebnisse einer Probe}

\subsubsection{Erfolg}

Wenn der Held genug Sterne würfelt, gelingt das Vorhaben. Die Spielleitung beschreibt, was passiert.

\textbf{Beispiel:} Lina will über einen Bach springen. Die Spielleitung sagt, dass dafür zwei Sterne nötig sind. Lina würfelt zwei Sterne. Sie springt sicher über den Bach und landet auf der anderen Seite.

Ein Erfolg bedeutet nicht immer, dass alles perfekt läuft. Aber das eigentliche Vorhaben gelingt.

Wenn ein Held besonders viele Sterne würfelt, kann die Spielleitung das Ergebnis besonders gut beschreiben.

\textbf{Beispiel:} Lina würfelt nicht nur zwei, sondern vier Sterne. Sie springt nicht nur über den Bach, sondern landet so geschickt, dass sie sofort weiterlaufen kann.

\subsubsection{Fehlschlag}

Wenn der Held zu wenige Sterne würfelt, gelingt das Vorhaben nicht oder nur teilweise.

Ein Fehlschlag soll die Geschichte nicht stoppen. Er soll zeigen, dass etwas schiefgeht und die Situation sich verändert.

Beispiele:
\begin{itemize}
  \item Der Held schafft es nicht rechtzeitig.
  \item Der Held macht Lärm.
  \item Der Held verliert etwas.
  \item Der Held braucht einen anderen Plan.
  \item Der Held gerät in Gefahr.
  \item Jemand bemerkt den Helden.
  \item Die Aufgabe dauert länger als gedacht.
\end{itemize}

\textbf{Beispiel:} Lina will über einen Bach springen. Sie braucht zwei Sterne, würfelt aber nur einen Stern. Sie schafft den Sprung nicht richtig. Vielleicht landet sie im Wasser. Vielleicht kommt sie auf der anderen Seite an, verliert dabei aber ihren Rucksack. Die Spielleitung entscheidet, was zur Situation passt.

Ein Fehlschlag bedeutet nicht: „Nichts passiert.“

Ein Fehlschlag bedeutet: „Etwas passiert anders als geplant.“

\subsubsection{Teilerfolg}

Manchmal passt weder ein voller Erfolg noch ein klarer Fehlschlag. Dann kann die Spielleitung einen Teilerfolg beschreiben.

Ein Teilerfolg bedeutet: Der Held schafft sein Vorhaben, aber es gibt ein Problem.

Beispiele:
\begin{itemize}
  \item Der Held kommt durch die Tür, aber macht dabei Lärm.
  \item Der Held findet die Spur, aber verliert viel Zeit.
  \item Der Held beruhigt das Tier, aber es läuft danach davon.
  \item Der Held klettert die Mauer hoch, aber lässt etwas fallen.
\end{itemize}

Teilerfolge sind besonders nützlich, wenn ein Fehlschlag die Geschichte unnötig blockieren würde.

\subsection{Vorteil und Nachteil}

Manchmal ist eine Probe leichter oder schwerer als gewöhnlich. Dann kann die Spielleitung Vorteil oder Nachteil geben.

Vorteil bedeutet: Die Umstände helfen dem Helden.

Nachteil bedeutet: Die Umstände machen es dem Helden schwerer.

Vorteil und Nachteil verändern nicht, was der Held versucht. Sie verändern nur, wie gut seine Chancen sind. Vorteil und Nachteil gelten nur für Proben von Helden. Sie gelten nicht für Schadenswürfe.

In seltenen Situationen kann ein Held gleichzeitig Vorteil und Nachteil erhalten. In diesem Fall heben sich Vorteil und Nachteil gegenseitig auf.

Nur Helden erhalten Vorteil und Nachteil. NPCs oder Gegner haben dies nie.

\subsubsection{Vorteil}

Ein Held kann Vorteil bekommen, wenn die Situation ihm hilft.

Beispiele:
\begin{itemize}
  \item Der Held hat gutes Werkzeug.
  \item Der Held hat viel Zeit.
  \item Der Held kennt den Ort.
  \item Der Held beschreibt eine besonders gute Idee.
  \item Die Umstände sind günstig.
\end{itemize}

Wenn ein Held Vorteil hat, darf er bei dieser Probe beliebige Würfel einmal kostenlos neu würfeln. Danach gilt das neue Ergebnis.

\subsubsection{Nachteil}

Ein Held kann Nachteil bekommen, wenn die Situation besonders schwierig ist.

Beispiele:
\begin{itemize}
  \item Der Held steht unter Zeitdruck.
  \item Der Held ist verletzt oder erschöpft.
  \item Es ist dunkel.
  \item Der Boden ist rutschig.
  \item Der Held hat schlechtes Werkzeug.
  \item Die Umstände sind gefährlich.
\end{itemize}

Wenn ein Held Nachteil hat, darf er bei dieser Probe keine Mondkristalle nutzen.

\subsection{Mondsymbole}

Der Heldenwürfel hat auf einigen Seiten Mondsymbole. Ein Mondsymbol zählt nicht als Stern. Wenn eine Würfelseite ein Mondsymbol und einen Stern zeigt, wird der Stern trotzdem normal gezählt.

Mondsymbole können verwendet werden, um besondere Fähigkeiten auszulösen. Jeder Held hat eine solche Fähigkeit von Anfang an auf seinem Charakterbogen: Er kann ein gewürfeltes Mondsymbol nutzen, um einen Mondkristall zu erhalten.

Ein Mondkristall ist ein Funken Schicksal, den die Helden einsetzen können, wenn sie Unterstützung brauchen.

Mondsymbole werden erst ausgewertet, nachdem alle kostenlosen Neuwürfe durch Vorteil abgeschlossen sind.

Wichtig: Sobald ein Würfel für einen Effekt genutzt wurde, ist er für diese Probe fest. Er kann in dieser Probe nicht mehr verändert oder neu geworfen werden.

\subsubsection{Mondkristalle}

Ein Held kann Mondkristalle jederzeit ausgeben, um eine Probe zu beeinflussen oder eine besondere Fähigkeit auszulösen.

Die Grundfähigkeit eines Mondkristalls ist: Der Held darf alle Würfel einer Farbe neu würfeln.

Zum Beispiel:
\begin{itemize}
  \item den weißen Heldenwürfel,
  \item alle grünen Attributswürfel,
  \item alle blauen Fertigkeitswürfel.
\end{itemize}

Der Spieler entscheidet, welche Würfelfarbe neu gewürfelt wird. Danach gilt das neue Ergebnis.

Ein Held darf in einer Probe mehrere Mondkristalle für die Grundfähigkeit ausgeben, wenn er genug Mondkristalle besitzt. Dadurch kann derselbe Würfel auch mehrfach neu geworfen werden, solange er noch nicht für einen Effekt genutzt wurde.

Neben der Grundfähigkeit gibt es Talente, die durch das Ausgeben von Mondkristallen ausgelöst werden können. Die genauen Regeln dazu stehen im Talente-Kapitel.

\subsection{Beispiel: Eine Probe am Spieltisch}

Kiera steht vor einer alten, halb eingestürzten Steinbrücke. Unter ihr rauscht ein kalter Fluss. Die Brücke ist wackelig, aber auf der anderen Seite hängt ein altes Seil, an dem Kiera sich festhalten kann.

Kieras Spieler sagt: „Ich laufe über die Brücke und halte mich dabei am Seil fest.“

Die Spielleitung entscheidet: Das ist möglich, aber gefährlich. Es wird eine Probe gemacht.

Die Spielleitung sagt: „Das passt zu Turnen. Die Probe ist schwierig. Du brauchst drei Sterne. Weil du das Seil nutzt, hast du Vorteil.“

Kieras Werte sind:
\begin{itemize}
  \item Gewandtheit 2,
  \item Kraft 1,
  \item Turnen 1.
\end{itemize}

Turnen gehört zu Gewandtheit und Kraft. Kieras Spieler nimmt daher:
\begin{itemize}
  \item den Heldenwürfel,
  \item 2 grüne Würfel für Gewandtheit,
  \item 1 grünen Würfel für Kraft,
  \item 1 blauen Würfel für Turnen.
\end{itemize}

Kiera würfelt zuerst. Der Heldenwürfel zeigt ein Mondsymbol, der blaue Turnenwürfel einen Stern und die drei grünen Würfel bleiben leer. Kiera hat damit zunächst nur einen Stern und braucht für die schwierige Probe drei.

Da sie Vorteil hat, würfelt sie die drei grünen Würfel einmal kostenlos neu. Einer davon zeigt nun einen Stern. Insgesamt hat Kiera jetzt zwei Sterne.

Kiera nutzt anschließend das Mondsymbol auf dem Heldenwürfel. Durch die Fähigkeit auf ihrem Charakterbogen erhält sie einen Mondkristall und gibt ihn sofort aus, um alle grünen Würfel erneut zu würfeln. Zwei grüne Würfel zeigen nun jeweils einen Stern.

Zusammen mit dem Stern auf dem blauen Turnenwürfel erreicht Kiera drei Sterne. Das reicht für die schwierige Probe.

Die Spielleitung beschreibt das Ergebnis: „Kiera setzt einen Fuß auf die wackelige Brücke. Ein Brett knackt unter ihr, und für einen Moment verliert sie fast das Gleichgewicht. Dann bekommt sie das alte Seil zu fassen, zieht sich nach vorne und nutzt den Funken Schicksal, den der Mond ihr schenkt. Sie findet wieder Halt und erreicht die andere Seite. Hinter ihr bricht ein Stück der Brücke in den Fluss. Kiera hat es geschafft.“
